\section{半年単位で再分類する — Model能力からExecution Responsibilityへ}
\label{analysis:h1-reclassification}

2025年前半の主要eventを月ごとのlaunch列として読むと、reasoning model、Agent、open-weight、multimodal、API/tool、safeguardが別々の流行に見える。
半年単位では、これらを\textbf{能力をどの層で実行可能にし、どこで制御するか}という責任軸へ再分類した方が変化を比較しやすい。
以下は本号で確認済みの一次資料を横断した編集上の整理であり、各vendorが共通architectureを主張しているという意味ではない。

\begin{center}
\small
\renewcommand{\arraystretch}{1.18}
\begin{tabularx}{\linewidth}{@{}p{0.19\linewidth}p{0.26\linewidth}X@{}}
\toprule
責任の層 & H1で前景化した対象 & 半年単位で見える変化 \\
\midrule
Reasoning & DeepSeek-R1、o3-mini、Claude 3.7、Gemini 2.5、o3/o4-mini &
reasoningは単独のmodel categoryではなく、research、tool use、coding、multimodal reasoningへ接続される基礎能力として扱われ始めた。
\cite{src-8f8555b73e1b,src-650b567698aa,src-a6c92f94a0af,src-421ee9af3f80,src-bf92f3d905c3} \\
Execution & Operator、Deep Research、Responses API、Claude Code、Codex、Jules &
model responseを返すだけでなく、browser、research、tool、repository上の作業を実行するsurfaceが独立した観測対象になった。
\cite{src-95b210494c88,src-8fec3ef3a97d,src-78593f4b387f,src-a6c92f94a0af,src-e7aef8f1b13b,src-8b595b787bad} \\
Delivery / Lifecycle & GPT-4.1、GPT-4.5 preview/retirement、Gemini 2.5 GA、API提供 &
「発表されたmodel」と「API/productとして利用可能なidentity」、さらにretirementを同じrelease eventへ潰さず追う必要が強まった。
\cite{src-31799a46d8d0,src-3fa7f4cd510b,src-53f15d483a2d} \\
Open Weight & Gemma 3、Llama 4、Qwen3 &
weight公開とAPI/product提供は別責任であり、open-weight family自体がfrontier競争の独立軸として再び存在感を増した。
\cite{src-1bb16d7254b1,src-05f285e9757b,src-3e4a3f3b3d83} \\
Multimodal / Media & Qwen2.5-VL、Qwen2.5-Omni、QVQ-Max、GPT-4o image、Veo 3等 &
multimodalityは入力対応だけでなく、visual reasoning、native image generation、video/image/music production workflowへ広がった。
\cite{src-dc21f366b638,src-de4ffd746276,src-8cb736b33646,src-877e7da88408,src-e68f8ce96685} \\
Control & Web Search、remote MCP、Llama Guard 4、interpretability &
execution surfaceの拡大と並行して、search/tool接続、safeguard、内部挙動の観測といったcontrol側の責任も独立して見え始めた。
\cite{src-e31c642d34b6,src-9257c39d4172,src-383476bdf62e,src-fdc10e3226e6} \\
\bottomrule
\end{tabularx}
\renewcommand{\arraystretch}{1.0}
\end{center}
\normalsize

\Needspace{0.38\textheight}
\section{Cross-month Comparison — 同じ責任を1月から6月まで追う}
\label{analysis:h1-cross-month}

H1をJan--Feb、Mar--Apr、May--Junの三段階に分けると、単発launchでは見えにくいresponsibility transitionが現れる。

\begin{center}
\footnotesize
\renewcommand{\arraystretch}{1.18}
\begin{tabularx}{\linewidth}{@{}p{0.16\linewidth}p{0.26\linewidth}p{0.26\linewidth}X@{}}
\toprule
責任 & Jan--Feb & Mar--Apr & May--Jun \\
\midrule
Reasoning &
DeepSeek-R1とo3-miniがreasoning modelを前面化し、Deep Researchがreasoningを長時間researchへ接続した。
\cite{src-8f8555b73e1b,src-650b567698aa,src-8fec3ef3a97d} &
Claude 3.7、Gemini 2.5、o3/o4-miniがreasoningをcoding・tool・multimodal capabilityと結び付けた。
\cite{src-a6c92f94a0af,src-421ee9af3f80,src-bf92f3d905c3} &
Gemini 2.5 Pro/Flash GAとClaude 4により、thinking/reasoningはpreviewだけでなく継続利用するmodel family能力として定着した。
\cite{src-53f15d483a2d,src-4ec03078d316} \\
Execution &
OperatorとDeep Researchがbrowser操作とmulti-step researchを別々のproduct surfaceとして示した。
\cite{src-95b210494c88,src-8fec3ef3a97d} &
Responses API / Agents SDKとClaude Codeが、tool・tracing・coding環境をmodel外側のexecution layerとして可視化した。
\cite{src-78593f4b387f,src-a6c92f94a0af} &
Codex、Jules、remote MCP/tool expansionが、repository作業と外部tool接続を継続的なAgent platformへ押し広げた。
\cite{src-e7aef8f1b13b,src-8b595b787bad,src-9257c39d4172} \\
Open / Multimodal &
Qwen2.5-VLとQwen2.5-Maxがopen modelとAPI availabilityを並行させた。
\cite{src-dc21f366b638,src-8c485019ad8d} &
Gemma 3、Llama 4、Qwen3とQwen2.5-Omni/QVQ-Maxが、open-weightとmultimodal reasoningを同時に拡張した。
\cite{src-1bb16d7254b1,src-05f285e9757b,src-3e4a3f3b3d83,src-de4ffd746276,src-8cb736b33646} &
GPT-4o imageとVeo 3等により、multimodalの焦点が理解から生成media workflowへ広がった。
\cite{src-877e7da88408,src-e68f8ce96685} \\
Control / Lifecycle &
research previewとAPI availabilityの区別が重要だった。 &
GPT-4.5 previewとGPT-4.1 API family、Llama Guard 4、interpretability研究が、lifecycleとcontrolをmodel releaseから分離して見せた。
\cite{src-3fa7f4cd510b,src-31799a46d8d0,src-383476bdf62e,src-fdc10e3226e6} &
Web Search API、remote MCP、Gemini 2.5 GAにより、外部接続とstable deliveryの責任が前景化した。
\cite{src-e31c642d34b6,src-9257c39d4172,src-53f15d483a2d} \\
\bottomrule
\end{tabularx}
\renewcommand{\arraystretch}{1.0}
\end{center}
\normalsize

\Needspace{0.36\textheight}
\section{Cross-layer Synthesis — H1で結合し始めた層}
\label{analysis:h1-cross-layer}

\paragraph{Reasoning $\times$ Execution}
DeepSeek-R1やo3系で前景化したreasoningは、Operator、Deep Research、Responses API、Claude Code、Codexへ接続されることで、
「より良い回答を生成する能力」だけでなく「複数stepを実行する能力」の一部になった。
\cite{src-8f8555b73e1b,src-bf92f3d905c3,src-95b210494c88,src-8fec3ef3a97d,src-78593f4b387f,src-e7aef8f1b13b}

\paragraph{Multimodal $\times$ Reasoning / Action}
Qwen2.5-VL、QVQ-Max、Gemini 2.5はvisual inputをreasoningへ取り込み、GPT-4o native image generationやVeo 3等は出力側をmedia productionへ拡張した。
理解と生成を同じ「multimodal」一語で扱うより、perception/reasoningとproduction/action spaceを分けて追う方がH1の変化を説明しやすい。
\cite{src-dc21f366b638,src-8cb736b33646,src-421ee9af3f80,src-877e7da88408,src-e68f8ce96685}

\paragraph{Open Weight $\times$ Delivery / Control}
Gemma 3、Llama 4、Qwen3のweight公開は、hosted APIとは異なるdelivery responsibilityを利用者側へ移す。
同じ4月末にはLlama API limited previewとLlama Guard 4も現れ、open-weight、hosted delivery、safeguardが同じecosystemでも別identityとして追跡すべきことが明確になった。
\cite{src-1bb16d7254b1,src-05f285e9757b,src-3e4a3f3b3d83,src-fca961b24baf,src-383476bdf62e}

\begin{claimboundary}[横断分析の境界]
本節の「責任」「層」「結合」は、本号で検証済みの独立したeventを同じ観測軸へ配置した編集上の推論である。
引用資料が共通のsystem architectureや因果関係を共同で主張しているわけではない。
H2で判明した後続結果をH1の事実として逆投影せず、2025年6月30日までに確認できるidentityとavailabilityを基準に整理している。
\end{claimboundary}
